Microgrid procurement links uncertain demand and photovoltaic (PV) generation to battery losses, discharge limits, and the cost of revising purchases. We develop schedules for the four tasks under explicit information and settlement rules. Source records are aligned by interval and forecast issue time; power is converted to energy, structural date labels are filled, and hourly PV forecasts are integrated without smoothing observations.

For Problem 1, a deterministic mixed-integer linear program jointly schedules purchases and battery flows. HiGHS returns a daily cost of \textbf{35,126.95 CNY}, \textbf{26.898\%} below the no-storage benchmark, with a reported zero optimality gap for that supplied instance.

For Problems 2-4, we use historical daily forecast errors to evaluate a shared purchase vector and a causal discharge-reserve vector. A threshold issued before execution can retain energy for later expensive deficits; setting it to 1,200 kWh recovers maximum deficit discharge. Four L-BFGS-B searches generate candidates for each decision, which are then explicitly rescored. Twenty-two continuous February-December 2025 paths compare the reserve and greedy controllers with the same search counts and per-search limits.

For Problem 2, the reserve policy costs \textbf{13,845,452.99 CNY} and saves \textbf{69,508.07 CNY} against its matched greedy control. Emergency expenditure nevertheless rises by \textbf{22,702.84 CNY}. For Problem 3, the selected corrected-PV rolling reserve policy costs \textbf{13,574,461.81 CNY}. Within the reserve family, state feedback with paid revisions saves CNY 208,115.60; adding later raw PV releases saves CNY 167,121.93. The selected corrected-PV reserve policy costs an additional CNY 6,206.07 relative to its matched greedy control, so additional storage flexibility does not assure lower realized cost.

For Problem 4, planning prices are separated from actual billed prices. The midnight branch selects the causal-OLS midnight reserve policy, costing \textbf{14,618,518.98 CNY}; the rolling branch selects the raw-PV/causal-OLS rolling reserve policy, costing \textbf{14,313,637.10 CNY}. Both use actual variable-price settlement and the full fivefold emergency tariff.

Independent reconstruction of source data, physical and fee replay, candidate checks, and isolated reproduction support the consistency of the implementation. The five workbooks contain the resulting schedules. Because we selected policies on a repeatedly inspected historical year and used finite nonsmooth searches, these results cannot establish future savings or global optimality.
